\باب{مقاطع اور کریمر کا قاعدہ}\شناخت{ضمیمہ_ح}
اعداد کی مستطیل  ترتیب :
\[
A=
\left[
\begin{array}{rrr}
2&1&3\\
1&0&-2
\end{array}
\right]
\]
\اصطلاح{قالب }\فرہنگ{قالب}\حاشیہب{matrix}\فرہنگ{matrix} کہلاتا ہے۔ اس قالب میں  دو  صف اور  تین  قطار ہیں لہٰذا   ہم \عددی{A} کو \اردوعدد{2}  با  \اردوعدد{3} قالب یا    \اردوعدد{2}  ضرب \اردوعدد{3} قالب  کہتے ہیں۔ ایک    \عددی{m} با \عددی{n} قالب     میں \عددی{m} صف اور \عددی{n}  قطار  ہوں گے، اور  قالب کے \عددی{i} ویں صف اور \عددی{j} ویں قطار میں موجود\اصطلاح{ رکن}\فرہنگ{قالب!رکن}\حاشیہب{element}\فرہنگ{matrix!element}  (عدد) \عددی{a_{ij}} سے ظاہر کیا جائے گا۔ یوں درج ذیل قالب
\[
A=
\left[
\begin{array}{rrr}
2&1&3\\
1&0&-2
\end{array}
\right]
\]
کے ارکان   درج ذیل  لکھے جائیں گے۔
\begin{align*}
a_{11}&=2,\quad a_{12}=1,\quad a_{13}=3,\\
a_{21}&=1,\quad a_{22}=0,\quad a_{23}=-2
\end{align*}

جس قالب میں صفوں  اور قطاروں کی تعداد برابر ہو \اصطلاح{مکعب قالب }\فرہنگ{قالب!مکعب}\حاشیہب{square matrix}\فرہنگ{matrix!square} کہلاتا ہے۔مکعب قالب  جس میں  صفوں کی تعداد اور قطاروں کی تعداد \عددی{n} ہو\اصطلاح{\عددی{n} رتبی قالب}\فرہنگ{قالب!رتبہ}\حاشیہب{matrix of order n}\فرہنگ{matrix!order} کہلاتا ہے۔

ہر مکعب قالب  \عددی{A} کے ساتھ ایک عدد   وابستہ کیا جاتا ہے  جو \عددی{A} کا \اصطلاح{مقطع }\فرہنگ{قالب!مقطع}\حاشیہب{determinant}\فرہنگ{matrix!determinant} کہلاتا  اور\قول{مقطع $A$} یا \عددی{|a_{ij}|} سے ظاہر کیا جاتا ہے۔مقطع کی قیمت قالب \عددی{A}  کے ارکان سے درج ذیل عمل کے ذریعہ  حاصل کیا جاتا ہے۔ یک رتبی \عددی{n=1} قالب اور  دو رتبی قالب \عددی{n=2} کے  مقطع کی تعریف درج ذیل ہے۔
\begin{align}
[a]\,\text{مقطع}\,&=\,a\\
\left[
\begin{array}{ll}
a_{11}& a_{12}\\
a_{21}&a_{22}
\end{array}
\right]\text{مقطع}\,
&=\,a_{11}a_{22}-a_{21}a_{12}
\end{align}
تین  رتبی قالب کے مقطع کی تعریف درج ذیل ہے
\begin{align}\label{مساوات_ضمیمہ_قالب_مقطع_تین}
A\,\text{مقطع}\,=\,\left[
\begin{array}{lll}
a_{11}&a_{12}&a_{13}\\
a_{21}&a_{22}&a_{23}\\
a_{31}&a_{32}&a_{33}
\end{array}
\right]\text{مقطع}\,=\,
\begin{tabular}{r}
\عددی{\pm a_{1i}a_{2j}a_{3k}} روپ کے تمام \\
علامت دار حاصل ضرب کا مجموعہ
\end{tabular}
\end{align}
جہاں \عددی{i}، \عددی{j}، \عددی{k}  سے مراد \عددی{1}، \عددی{2}، \عددی{3} کی کسی ترتیب سے مرتب اجتماع ہے۔ ایسے \عددی{3!=6} مرتب اجتماعات ہیں لہٰذا مجموعے  میں چھ جزو ہوں گے۔ جب مرتب اجتماعات کا اشاریہ  جفت ہو، علامت مثبت  ہو گی اور جب طاق ہو علامت منفی ہو گی۔

\ابتدا{تعریف}\موٹا{مرتب اجتماع کا اشاریہ}\\
اعداد \عددی{1,2,3,\cdots,n} کی  دی گئی مرتب اجتماع  کو  \عددی{i_1,i_2,i_3,\cdots,i_n}  سے ظاہر کریں۔ اس نظم میں \عددی{i_1} کے بعد آنے والے   چند اعداد \عددی{i_1} سے چھوٹے ہو سکتے ہیں۔ایسے اعداد کی تعداد  کو  نظم میں \عددی{i_1}  سے متعلق \اصطلاح{تعکیس کی تعداد }\فرہنگ{تعداد!تعکیس}\حاشیہب{number of inversions}\فرہنگ{number!inversions} کہتے ہیں۔ اسی طرح نظم میں تمام   دیگر \عددی{i}  سے متعلق  تعکیس کی تعداد ہو گی، جو نظم میں  اس مخصوص \عددی{i} کے بعد  ان اعداد کی تعداد ہو گی جو اس مخصوص عدد سے چھوٹے ہوں۔ تمام اشاریوں کے تعکیس کی تعداد کا مجموعہ  مرتب اجتماع کا \اصطلاح{اشاریہ }\فرہنگ{مرتب اجتماع!اشاریہ}\حاشیہب{index}\فرہنگ{permutation!index} کہلاتا ہے۔
\انتہا{تعریف}

\ابتدا{مثال}
عدد \عددی{n=5}  کی ایک مرتب اجتماع  ذیل ہے
\[5\hspace{1em} 3 \hspace{1em}  1\hspace{1em} 2 \hspace{1em} 2\]
جس میں پہلا  رکن \عددی{5} ہے جس کے بعد آنے والے  چاروں  ارکان  اس سے چھوٹے ہیں لہٰذا  \عددی{5} سے متعلق تعکیس کی تعداد  \عددی{4}  ہے؛ دوسرا رکن \عددی{3} ہے  اور اس کے بعد آنے والے دو  ارکان      اس سے چھوٹے ہیں لہٰذا  \عددی{3} سے متعلق تعکیس کی تعداد  \عددی{2}  ہے؛ اگلا  رکن  \عددی{1} ہے  جس کے بعد ایسا کوئی رکن نہیں آتا جو ایک  سے چھوٹا ہو لہٰذا \عددی{1} سے متعلق تعکیس کی تعداد  \عددی{0}  ہے؛ اگلے دونوں ارکان  سے متعلق  انفرادی تعکیس کی تعداد \عددی{0}  ہے؛ یوں اشاریہ \عددی{4+2=6} ہو گا۔
\انتہا{مثال}


اگلا جدول \عددی{1,2,3} کی مرتب اجتماعات ، ہر مرتب اجتماع کا اشاریہ، اور  مساوات \حوالہ{مساوات_ضمیمہ_قالب_مقطع_تین}  کے مقطع  میں علامت دار حاصل ضرب دیتا ہے۔ 
\begin{align}
\begin{tabular}{ccc}
مرتب اجتماع&اشاریہ&علامت دار حاصل ضرب\\
\midrule
\(1 \hspace{1em} 2  \hspace{1em} 3\)&\(0\)&\(+a_{11}a_{22}a_{33}\) \\
\(1 \hspace{1em} 3  \hspace{1em} 2\)&\(1\)&\(-a_{11}a_{23}a_{32}\) \\
\(2 \hspace{1em} 1  \hspace{1em} 3\)&\(1\)&\(-a_{12}a_{21}a_{33}\) \\
\(2 \hspace{1em} 3  \hspace{1em} 1\)&\(2\)&\(+a_{12}a_{23}a_{31}\) \\
\(3 \hspace{1em} 1  \hspace{1em} 2\)&\(2\)&\(+a_{13}a_{21}a_{32}\) \\
\(3 \hspace{1em} 2  \hspace{1em} 1\)&\(3\)&\(-a_{13}a_{22}a_{31}\) 
\end{tabular}
\end{align}

ان چھ حاصل ضرب کا مجموعہ ذیل ہو گا۔
\begin{align}
a_{11}(a_{22}a_{33}-a_{23}a_{32})&-a_{12}(a_{21}a_{33}-a_{23}a_{31})+a_{13}(a_{21}a_{32}-a_{22}a_{31})\notag\\
&=a_{11}
\left|
\begin{array}{ll}
a_{22}&a_{23}\\
a_{32}&a_{33}
\end{array}\right|
-a_{12}
\left|
\begin{array}{ll}
a_{21}&a_{23}\\
a_{31}&a_{33}
\end{array}\right|
+a_{13}
\left|
\begin{array}{ll}
a_{21}&a_{22}\\
a_{31}&a_{32}
\end{array}\right| \notag\\
&=
\left|
\begin{array}{lll}
a_{11}& a_{12}&a_{13}\\
a_{21}& a_{22}&a_{23}\\
a_{31}& a_{32}&a_{33}
\end{array}\right|
\end{align}
درج ذیل کلیہ \عددی{3}  با \عددی{3} مقطع کے حساب کو تین عدد \عددی{2} با \عددی{2} مقطع کے حساب میں تبدیل کرتا ہے۔
\begin{align}\label{مساوات_ضمیمہ_قالب_مقطع_سادہ}
\left|
\begin{array}{lll}
a_{11}& a_{12}&a_{13}\\
a_{21}& a_{22}&a_{23}\\
a_{31}& a_{32}&a_{33}
\end{array}\right|=
a_{11}
\left|
\begin{array}{ll}
a_{22}&a_{23}\\
a_{32}&a_{33}
\end{array}\right|
-a_{12}
\left|
\begin{array}{ll}
a_{21}&a_{23}\\
a_{31}&a_{33}
\end{array}\right|
+a_{13}
\left|
\begin{array}{ll}
a_{21}&a_{22}\\
a_{31}&a_{32}
\end{array}\right|
\end{align}

\جزوحصہء{اصاغر  اور ہمزاد اجزائے   ضربی}
مساوات \حوالہ{مساوات_ضمیمہ_قالب_مقطع_سادہ} کے دائیں ہاتھ دو رتبی مقاطع  کو ان ارکان کے   \اصطلاح{اصاغر }\فرہنگ{اصاغر}\حاشیہب{minors}\فرہنگ{minors} (صغیر مقطع کی مختصر اصطلاح)    کہتے ہیں جن سے انہیں ضرب کیا گیا ہے۔یوں\\
\[\left|
\begin{array}{ll}
a_{21}&a_{23}\\
a_{31}&a_{33}
\end{array}
\right|\hspace{1em}  \text{\RL{اور \عددی{a_{12}} کا اصغر}}\quad
\left|
\begin{array}{ll}
a_{22}&a_{23}\\
a_{32}&a_{33}
\end{array}
\right|\hspace{1em}  \text{\RL{\عددی{a_{11}} کا اصغر}}
\]
ہو گا، وغیرہ۔ قالب \عددی{A}  سے  وہ صف اور  قطار خارج کرنے کے بعد ، جس میں رکن \عددی{a_{ij}} پایا جاتا ہو، حاصل قالب  کا مقطع \عددی{a_{ij}} کا اصغر ہو گا۔
%==================

\(\begin{vmatrix}
a_{11}& \tikzmark{start}a_{12}&a_{13}\\
\tikzmark{startL}a_{21}& a_{22}&a_{23}\tikzmark{endL}\\
a_{31}& a_{32}\tikzmark{end}&a_{33}
\end{vmatrix}\)\quad
سے   \عددی{a_{22}}  کا  اصغر    \quad
\(\begin{vmatrix}
a_{11}&a_{13}\\
a_{31}&a_{33}
\end{vmatrix}\)\quad
حاصل ہوتا ہے۔
\begin{tikzpicture}[remember picture, overlay]
  \draw[thick] (start.north) -- (end.south);
  \draw[thick] (startL.east) -- (endL.west);
\end{tikzpicture}

%================================

\(\begin{vmatrix}
a_{11}& a_{12}&\tikzmark{start}a_{13}\\
\tikzmark{startL}a_{21}& a_{22}&a_{23}\tikzmark{endL}\\
a_{31}& a_{32}&a_{33}\tikzmark{end}
\end{vmatrix}\)\quad
سے   \عددی{a_{23}}  کا  اصغر    \quad
\(\begin{vmatrix}
a_{11}&a_{12}\\
a_{31}&a_{32}
\end{vmatrix}\)\quad
حاصل ہوتا ہے۔
\begin{tikzpicture}[remember picture, overlay]
  \draw[thick] (start.north) -- (end.south);
  \draw[thick] (startL.east) -- (endL.west);
\end{tikzpicture}

%=======================
 رکن \عددی{a_{ij}} کا ہمزاد جزو  ضربی  \عددی{A_{ij}}  حاصل کرنے کے لئے \عددی{a_{ij}} کے اصغر کو \عددی{(-1)^{i+j}} سے ضرب دیں۔یوں درج ذیل ہوں گے۔
\begin{align*}
A_{22}&=(-1)^{2+2} 
\begin{vmatrix}
a_{11}&a_{13}\\
a_{31}&a_{33}
\end{vmatrix}\,=\,
\begin{vmatrix}
a_{11}&a_{13}\\
a_{31}&a_{33}
\end{vmatrix}\\
A_{23}&=(-1)^{2+3} 
\begin{vmatrix}
a_{11}&a_{12}\\
a_{31}&a_{32}
\end{vmatrix}\,=\,
-\begin{vmatrix}
a_{11}&a_{12}\\
a_{31}&a_{32}
\end{vmatrix}
\end{align*}
جزو ضربی \عددی{(-1)^{i+j}} اصغر کی علامت  اس صورت تبدیل کرتا ہے جب \عددی{i+j} طاق ہو۔درج ذیل  نقشہ اس کی وضاحت کرتا  ہے۔
\[\begin{matrix}
+&-&+\\
-&+&-\\
+&-&+
\end{matrix}\]

اوپر بائیں کونے میں \عددی{i=1}، \عددی{j=1} ہے لہٰذا \عددی{(-1)^{1+1}=+1}  ہے۔اسی صف یا اسی  قطار میں رہتے ہوئے اگلے خانے میں \عددی{i}  یا \عددی{j} کی قیمت بڑھ کر  \عددی{2} ہو جاتی ہے، تاہم \عددی{i} اور \عددی{j} دونوں کی قیمت تبدیل نہیں ہوتی، اور یوں \عددی{(-1)^{1+2}=-1} حاصل ہو گا۔آپ دیکھ سکتے ہیں کہ  کسی ایک صف یا کسی  ایک قطار میں رہ کر  چلتے ہوئے ہر قدم پر   علامت  \عددی{+} سے \عددی{-} یا \عددی{-} سے \عددی{+} ہو گی۔

ہم مساوات \حوالہ{مساوات_ضمیمہ_قالب_مقطع_سادہ}کو ہمزاد اجزائے ضربی کے روپ میں درج لکھ سکتے ہیں۔
\begin{align}
A\,\text{}\,=\, a_{11}A_{11}+a_{12}A_{12}+a_{13}A_{13}
\end{align}
